\subsection{Función estructural de \texorpdfstring{\(K\) y \(\alpha\)}{K y alfa} en la incidencia excepcional}

El enlace entre la determinación numérica y la realización excepcional
admite una descripción operacional. En la lectura aritmética, \(K\)
actúa en la normalización de los registros regionales. En la lectura de
incidencia, sus componentes seleccionan el soporte
\[
 B_K=\{j:K_j\geq729\}=\{4,6,9,10\}.
\]
El mismo registro ordenado participa, por tanto, en dos funciones
especificadas: una composición posicional y una selección de incidencia.
Su realización racional conserva el registro cuando se fijan base,
longitud y origen, mientras el paso al soporte retiene la información
necesaria para la bandera combinatoria.

El registro firmado previo a la normalización de \(\alpha\) aporta
la hexada \(H_\alpha^-\). Las intersecciones de esta hexada con los
soportes regionales y con la incidencia de APP determinan el punto
marcado. Ese punto individualiza el vector de norma \(54\) empleado
en el vecino del retículo. Así, la orientación que interviene en la
normalización aritmética vuelve a actuar en la selección geométrica.
La relación no requiere deducir un retículo a partir de una aproximación
decimal: utiliza los datos firmados que preceden a su evaluación.

La realización excepcional conserva sucesivamente tipos diferentes de
relación. El código conserva combinaciones lineales; el diseño conserva
incidencias entre soportes; el pegado conserva clases discriminantes; el
retículo conserva su forma bilineal; el álgebra de vértices conserva la
graduación y los productos correspondientes. Sus grupos de automorfismos
se definen respecto de esas estructuras. La aparición del Monstruo al
final de la construcción expresa una propiedad del álgebra
\(V^\natural\) obtenida mediante el procedimiento citado, y Moonshine
relaciona las trazas graduadas de esa acción con funciones modulares.

Esta jerarquía proporciona un significado matemático concreto a la
función de enlace atribuida a \(K\) y \(\alpha\). El primero es un
registro integral reversible con realizaciones posicionales e
incidenciales; la segunda es la coordenada de compatibilidad aritmética
cuyo registro orientado interviene también en la incidencia marcada.
Las definiciones adquieren contenido por esas operaciones y por las
identidades expuestas en este capítulo. El desarrollo de la incidencia
es parte sustantiva del argumento del artículo: muestra qué relaciones
subsisten cuando la construcción deja de evaluarse como una coordenada
real y se realiza como estructura algebraica.
